% !TEX root = task-2-9-2.tex
% Editable task bank. Append new BankTask blocks at the end of the document.
\documentclass[12pt,a4paper]{article}
\usepackage{physics-textbook}
\usepackage{graphicx,needspace}
\usetikzlibrary{decorations.pathmorphing}
\newif\ifShowSolutions
\ifdefined\TutorEdition
  \ShowSolutionstrue
\else
  \ShowSolutionsfalse % Change to \ShowSolutionstrue to display solutions.
\fi
\newcounter{banktask}
\newsavebox{\BankTaskBox}
\newcommand{\BankTaskContent}[3]{%
  \noindent\textbf{\thebanktask.}\quad #2\par
  \ifShowSolutions
    \Needspace{50mm}% Keep the solution heading with its opening text or diagram.
    #3\par
    \ifdefined\TutorEdition
      \ifcsname TutorSource#1\endcsname
        \nopagebreak\smallskip{\small\csname TutorSource#1\endcsname\par}%
      \fi
    \fi
  \fi
}
\newcommand{\BankTask}[3]{%
  \refstepcounter{banktask}\par\medskip
  \sbox{\BankTaskBox}{\begin{minipage}{\linewidth}%
    \BankTaskContent{#1}{#2}{#3}\end{minipage}}%
  \ifdim\dimexpr\ht\BankTaskBox+\dp\BankTaskBox\relax>\dimexpr\textheight-20mm\relax
    \Needspace{35mm}\BankTaskContent{#1}{#2}{#3}%
  \else
    \Needspace{\dimexpr\ht\BankTaskBox+\dp\BankTaskBox\relax}%
    \noindent\usebox{\BankTaskBox}\par
  \fi
  \medskip
}
\begin{document}
\BookSubsection{2.9.2}{Уравнение теплового баланса}


% Source sdamgia-12859: https://phys-ege.sdamgia.ru/problem?id=12859
\BankTask{sdamgia-12859}{
В калориметр налит $1\,\text{л}$ воды при температуре $0^\circ\text{C}$. В этот калориметр последовательно выливают $50$ одинаковых мензурок воды, нагретой до температуры +$50^\circ\text{C}$. Объем мензурки $20\,\text{см}^{3}$. Потерями теплоты и теплоемкостью калориметра можно пренебречь. Какая температура (в °С) установится в калориметре?
}{
\smallskip\noindent\textbf{Решение}\par
Объем добавленной воды $50\cdot 20\,\text{см}^{3}$ = $1000\,\text{см}^{3}$ = $1\,\text{л}$ совпадает с объемом исходной воды, значит, температура установится средней между их температурами, то есть $25^\circ\text{C}$.

Ответ: $25$.
}

% Source sdamgia-19663: https://phys-ege.sdamgia.ru/problem?id=19663
\BankTask{sdamgia-19663}{
В калориметр залили три порции воды массами $200\,\text{г}$, $300\,\text{г}$ и $500\,\text{г}$, которые имели температуры $20^\circ\text{C}$, $40^\circ\text{C}$ и $60^\circ\text{C}$ соответственно. Теплообмен воды с окружающими телами пренебрежимо мал. Какой будет температура воды в калориметре после установления теплового равновесия? (В градусах Цельсия.)
}{
\smallskip\noindent\textbf{Решение}\par
Уравнение теплообмена: $Q_{\text{получ}}=Q_1 + Q_2 + Q_3=0.$ Теплоту вычислим по формуле $Q = cm_x ( t - t_x ) ,$ откуда $Q_{\text{получ}} = cm_1 ( t - t_1 ) + cm_2 ( t - t_2 ) + cm_3 ( t - t_3 ) = 0.$

Раскрываем скобки, приводим подобные и получаем выражение для $t$:

$t = \frac{m_1t_1 + m_2t_2 + m_3t_3}{m_1 + m_2 + m_3} = \frac{0{,}2 \cdot 20 + 0{,}3 \cdot 40 + 0{,}5 \cdot 60}{0{,}2 + 0{,}3 + 0{,}5} = 46 ^\circ\text{C}.$

Ответ: $46$.
}

% Source sdamgia-25242: https://phys-ege.sdamgia.ru/problem?id=25242
\BankTask{sdamgia-25242}{
В медный калориметр массой $500\,\text{г}$, находившийся при температуре $20^\circ\text{C}$, налили $200\,\text{г}$ воды, температура которой была равна $50^\circ\text{C}$. Через некоторое время, когда между водой и калориметром установилось тепловое равновесие, температура воды в нем оказалась равной $40^\circ\text{C}$. Какое количество теплоты было отдано водой и калориметром в окружающую среду? (В джоулях.)
}{
\smallskip\noindent\textbf{Решение}\par
Теплота, полученная калориметром, равна $Q_1=m_1c_1 ( t - t_1 ) =0{,}5 \cdot 380 \cdot ( 40 - 20 ) =3800\text{Дж}.$

Теплота, отданная водой, равна $Q_2=m_2c_2 ( t - t_2 ) =0{,}2 \cdot 4200 \cdot ( 40 - 50 ) = - 8400\text{Дж}.$

Теплота, отданная в окружающую среду, равна $Q=|Q_2| - Q_1=8400 - 3800=4600\text{Дж}.$

Ответ: $4600$.
}

% Source sdamgia-35343: https://phys-ege.sdamgia.ru/problem?id=35343
\BankTask{sdamgia-35343}{
В стакан калориметра, содержащий $250\,\text{г}$ воды, опустили кусок льда массой $140\,\text{г}$, имевшего температуру $0^\circ\text{C}$. После того как наступило тепловое равновесие, весь лед растаял, и температура воды стала равной $0^\circ\text{C}$. Определите начальную температуру воды. Теплоёмкостью калориметра
и теплообменом с окружающей средой пренебречь.
}{
\smallskip\noindent\textbf{Решение}\par
Поскольку потери по условию отсутствуют, то все количество теплоты, отданное водой при охлаждении от начальной температуры $t_1$ до температуры плавления
\[Q_1=m_{\text{в}}c(t_{\text{пл}}-t_1),\]
пошло на плавление льда
\[Q_2=m_{\text{л}}\lambda.\]
Запишем уравнение теплового баланса
\[Q_1+Q_2=0,\]
или
\[cm_{\text{в}}(t_1-t_2)=\lambda m_{\text{л}},\]
откуда
\[t_1=\frac{\lambda m_{\text{л}}}{cm_{\text{в}}}+t_2.\]
Подставив числовые значения, получим искомую величину:
\[t_1=\frac{3{,}3\cdot 10^5\cdot0{,}14}{4200\cdot0{,}25}+0=44^\circ\text{C}.\]
Ответ: $44^\circ\text{C}$.
}

% Source sdamgia-41579: https://phys-ege.sdamgia.ru/problem?id=41579
\BankTask{sdamgia-41579}{
Кусок льда опустили в термос с водой массой $1100\,\text{г}$, имеющей начальную температуру $15 ^\circ\text{C} .$ Начальная температура льда равна $0 ^\circ\text{C} .$  При достижении теплового равновесия в воде остаётся плавать кусочек льда. Какая масса льда растаяла в процессе перехода к тепловому равновесию? Теплоёмкостью термоса можно пренебречь.
}{
\smallskip\noindent\textbf{Решение}\par
Так как начальная температура льда равна температуре плавления, а при достижении теплового равновесия в воде остаются плавать кусочки льда, то делаем вывод, что установившаяся температура в сосуде равна $0$ \ensuremath{{}^\circ}C. При охлаждении воды выделилось количество теплоты $Q_1=mc ( t_{\,\text{пл}} - t ) ,$ на плавление льда потребовалось количество теплоты $Q_2=\lambda m_1.$
Из уравнения теплового баланса $Q_1 + Q_2=0$ находим массу растаявшего льда

\[m_1= \frac{mc ( t - t_{\,\text{пл}} )}{\lambda} = \frac{4200 \cdot 1{,}1 \cdot 15}{3{,}3 \cdot 10 ^{5}} =0{,}21 \,\text{кг}=210\,\text{г}.\]

Ответ: $210\,\text{г}$.
}

% Source sdamgia-3805: https://phys-ege.sdamgia.ru/problem?id=3805
\BankTask{sdamgia-3805}{
Кусок льда, имеющий температуру $0^\circ\text{C}$, помещен в калориметр с электронагревателем. Чтобы превратить этот лед в воду с температурой $10^\circ\text{C}$, требуется количество теплоты $200\,\text{кДж}$. Какая температура установится внутри калориметра, если лед получит от нагревателя количество теплоты $120\,\text{кДж}$? Теплоемкостью калориметра и теплообменом с внешней средой пренебречь. (В градусах Цельсия.)
}{
\smallskip\noindent\textbf{Решение}\par
Тогда для того, чтобы растопить лед массой $m$ и нагреть его до $10^\circ\text{C}$ требуется количество теплоты:

$Q=\lambda m + cm ( t - t_{\text{пл}} ) .$

По условию на это потребовалось $200\,\text{кДж}$. Следовательно, масса льда приблизительно равна:

$m= \frac{Q}{\lambda + c ( t - t_{\text{пл}} )} = \frac{200\text{кДж}}{3{,}3 \cdot 10 ^{5} \text{Дж}/\text{кг} + 4200\text{Дж}/ ( \text{кг} \cdot ^\circ\text{C} ) \cdot 10 ^\circ\text{C}} \approx 0{,}54 \text{кг}.$

Рассмотрим, сколько нужна тепла, чтобы растопить $0{,}54\,\text{кг}$ льда:

$Q_1=\lambda m=3{,}3 \cdot 10 ^{5} \text{Дж}/\text{кг} \cdot 0{,}54\text{кг} \approx 178\text{кДж}.$

Таким образом, если лед получит от нагревателя $120\,\text{кДж}$, то он не успеет растаять полностью, а значит, температура содержимого калориметра будет равна $0^\circ\text{C}$.

Ответ: $0^\circ\text{C}$.
}

% Source sdamgia-9290: https://phys-ege.sdamgia.ru/problem?id=9290
\BankTask{sdamgia-9290}{
В калориметр, содержащий $200\,\text{г}$ воды при температуре $15^\circ\text{C}$, добавили $20\,\text{г}$ мокрого снега. Температура в калориметре стала равна $10^\circ\text{C}$. Сколько воды было в снеге?
}{
\smallskip\noindent\textbf{Решение}\par
Обозначим за $m_{\text{в}}$ массу воды в снеге, тогда масса льда $m - m_{\text{в}},$ где $m$ — масса всего снега.

Количество теплоты, полученное льдом при его таянии при $0 ^\circ\text{C}:$ 

$Q_1=\lambda ( m - m_{\text{в}} ) .$

Количество теплоты, полученное при нагреве воды из всего снега до конечной температуры $t=10 ^\circ\text{C}$: 

$Q_2=cm ( t - 0 ) .$

Количество теплоты, отданное водой в калориметре при охлаждении ее от $t_1=15 ^\circ\text{C}$ до $t=10 ^\circ\text{C}$: 

$Q=cM ( t_1 - t ) .$ 

Уравнение теплового баланса: $Q=Q_1 + Q_2.$ 

Объединяя все выражения, получаем: 

$m_{\text{в}}=m - \frac{c ( M ( t_1 - t ) - mt )}{\lambda} \approx 10\text{г}.$

Ответ: $10$.
}

% Source sdamgia-20038: https://phys-ege.sdamgia.ru/problem?id=20038
\BankTask{sdamgia-20038}{
В калориметре находятся в тепловом равновесии вода и лед. После опускания в калориметр болта, имеющего массу $165\,\text{г}$ и температуру –$40$ \ensuremath{{}^\circ}С, $20$\% воды превратилось в лед. Удельная теплоемкость материала болта равна $500\,\text{Дж}$/(кг · К). Какая масса воды первоначально находилась в калориметре? Теплоемкостью калориметра пренебречь.
}{
\smallskip\noindent\textbf{Решение}\par
Так как вода и лед находятся в тепловом равновесии, то и до опускания болта, и после его нагревания температура в сосуде $t_{0}$  =  $0$ \ensuremath{{}^\circ}С. Согласно уравнению теплового баланса количество теплоты, выделившееся при замерзании воды, было затрачено на нагревание болта: $0{,}2m \cdot r = cm_1 ( t_0 - t ) ,$ где m  — масса воды в сосуде, $m_{1}$  — масса болта, с  — удельная теплоемкость болта, r  — удельная теплота плавления льда, t  — начальная температура болта. Получим:

\[m= \frac{cm_1 ( t_0 - t )}{0{,}2r} = \frac{500 \cdot 0{,}165 \cdot 40}{0{,}2 \cdot 3{,}3 \cdot 10 ^{5}} =0{,}05 \,\text{кг}.\]

Ответ: m  =  $0{,}05\,\text{кг}$.
}

% Source sdamgia-10963: https://phys-ege.sdamgia.ru/problem?id=10963
\BankTask{sdamgia-10963}{
В вертикальный теплоизолированный стакан калориметра объемом $200\,\text{см}^{3}$ налили до краев воду при температуре $t_{1}=20^\circ\text{C}$, а затем опустили туда кусок алюминия массой $m=270\,\text{г}$, находящийся при температуре $t_{2}$ = -$100^\circ\text{C}$. Какой объем льда окажется в стакане после установления теплового равновесия? Теплоемкостью стакана и поверхностным натяжением воды можно пренебречь. Плотность льда $0{,}9\,\text{г/см}^{3}$.
}{
\smallskip\noindent\textbf{Решение}\par
$1$. Выясним, какое количество теплоты необходимо для нагревания куска алюминия от температуры $t_{2}$ до $t_{0}=0^\circ\text{C}$: 

$Q_1=C_{\text{ал}}m ( t_0 - t_2 ) =900 \cdot 0{,}27 \cdot 100=24300\text{Дж}.$

$2$. Объем куска алюминия равен $V_2= \frac{m}{\rho_{\text{ал}}} = \frac{270}{2{,}7} =100\text{см} ^{3} ,$ и после его погружения в стакан часть воды вытечет, ее объем уменьшится на величину $V_{2}$ и станет равным $V_{1}=100\,\text{см}^{3}$, а ее масса будет равна:

$m_1=\rho_{\text{в}}V_1=1 \cdot 100=100\text{г}.$

$3$. Найдем теперь, какое количество теплоты может выделиться при охлаждении массы $m_1$ воды от температуры $t_{1}$ до $t_{0}=0^\circ\text{C}$: 

$Q_2=C_{\text{в}}m_1 ( t_1 - t_0 ) =4200 \cdot 0{,}1 \cdot 20=8400\text{Дж}.$

$4$. Поскольку $Q_{1}$ > $Q_{2}$, то часть воды начнет замерзать при $t_{0}=0^\circ\text{C}$ и ее теплота кристаллизации пойдет на нагревание алюминия. Всего замерзнет масса воды, равная $m_3= \frac{Q_1 - Q_2}{\lambda} = \frac{15900}{330} \approx48{,}18\text{г},$ которая займет объем:

$V_3= \frac{m_3}{\rho_{\text{л}}} = \frac{48{,}18}{0{,}9} \approx53{,}5\text{см} ^{3}$ 

в виде льда, примерзшего к куску алюминия.

Ответ: $V_{\text{л}}\approx53{,}5\text{см} ^{3} .$
}

% Source sdamgia-43507: https://phys-ege.sdamgia.ru/problem?id=43507
\BankTask{sdamgia-43507}{
В двух сосудах находится одинаковое количество воды при одинаковой температуре. В первый сосуд добавляют $0{,}42\,\text{кг}$ воды при температуре $100 ^\circ\text{C},$ а во второй — $0{,}42\,\text{кг}$ пара при температуре $100 ^\circ\text{C}.$ Конечная температура воды в сосудах после установления теплового равновесия отличается на $50 ^\circ\text{C}.$ Найдите начальную массу воды в каждом сосуде. Теплообменом с окружающей средой и теплоемкостью сосудов пренебречь.
}{
\smallskip\noindent\textbf{Решение}\par
В первом случае происходит теплообмен между двумя порциями воды, при котором холодная вода нагревается, получив количество теплоты $Q_1=cm ( t_1 - t_0 ) ,$ где $t_{1}$ — температура смеси в первом сосуде. Горячая вода охлаждается, отдавая количество теплоты $Q_2=cM ( t_1 - t_{\text{к}} ) .$ При отсутствии тепловых потерь уравнение теплового баланса:

$cm ( t_1 - t_0 ) + cM ( t_1 - t_{\text{к}} ) =0.$

Во втором сосуде пар конденсируется при температуре кипения, отдавая количество теплоты $Q_1= - LM,$ полученная вода охлаждается до установления температуры $t_{2}$, отдавая количество теплоты $Q_2=cM ( t_2 - t_{\text{к}} ) .$ Холодная вода нагревается, получив количество теплоты $Q_3=cm ( t_2 - t_0 ) .$ Тогда уравнение теплового баланса для второго сосуда:

$- LM + cm ( t_2 - t_0 ) + cM ( t_2 - t_{\text{к}} ) =0.$

Учитывая, что $\Delta t=t_2 - t_1,$ находим массу холодной воды:

$m= \frac{M ( L - c\Delta t )}{c\Delta t} = \frac{0{,}42 \cdot ( 2{,}3 \cdot 10 ^{6} - 4200 \cdot 50 )}{4200 \cdot 50} =4{,}18\text{кг}.$

Ответ: $4{,}18\,\text{кг}$.
}


% Source 70D177: https://ege.fipi.ru/bank/questions.php?search=1&proj=BA1F39653304A5B041B656915DC36B38&qid=70D177
\BankTask{70D177}{
В калориметр со льдом, имеющим температуру \(0\,^{\circ}\mathrm{C}\), добавили воду при комнатной температуре. Как изменятся в результате установления теплового равновесия удельная теплоёмкость воды и масса льда в калориметре?

Для каждой величины определите соответствующий характер изменения:
\begin{enumerate}
\item увеличится
\item уменьшится
\item не изменится
\end{enumerate}

\TaskChangeTable{Удельная теплоёмкость воды}{Масса льда}
}{
\smallskip\noindent\textbf{Решение от GPT}\par
Удельная теплоёмкость воды в принятой модели постоянна: \(c=4200\,\text{Дж}/(\text{кг}\cdot\text{К})\), поэтому не изменится (3). Вода при охлаждении отдаёт теплоту, которая расходуется на плавление льда: \(Q=\lambda\Delta m\). Масса льда уменьшится (2).

\(\boxed{32}\).
}

% Source 38577: https://phys-ege.sdamgia.ru/problem?id=38577
\BankTask{38577}{
В кастрюле находится $0{,}5\,\text{кг}$ воды при температуре $10 ^\circ\text{C}$. Сколько потребуется времени, чтобы при помощи кипятильника с постоянной потребляемой мощностью $400\,\text{Вт}$ превратить в пар $15\%$ воды из кастрюли? Потерями тепла и теплоёмкостью кастрюли пренебречь.
}{
\smallskip\noindent\textbf{Решение}\par
Вода в кастрюле нагревается до температуры кипения $100^\circ\text{C}$, получив количество теплоты
\[Q_1=mc(t_{\text{к}}-t),\]
а затем частично превращается в пар, получив количество теплоты
\[Q_2=0{,}15mL.\]
Так как потерь тепла не происходило, то это количество теплоты равно количеству теплоты, полученному от кипятильника. Из формулы тепловой мощности $P=\dfrac{Q}{\tau}$ при объединении формул получаем:
\[P=\frac{m\bigl(c(t_{\text{к}}-t)+0{,}15L\bigr)}{\tau},\]
откуда находим необходимое время:
\[\begin{aligned}
\tau&=\frac{m\bigl(c(t_{\text{к}}-t)+0{,}15L\bigr)}{P}\\
&=\frac{0{,}5\bigl(4{,}2\cdot10^3(100-10)+0{,}15\cdot2{,}3\cdot10^6\bigr)}{400}\approx904\,\text{с}.
\end{aligned}\]
Ответ: $904\,\text{с}$.
}

% Source 3435: https://phys-ege.sdamgia.ru/problem?id=3435
\BankTask{3435}{
В теплоизолированный сосуд с большим количеством льда при температуре $t_1=0 ^\circ\,\text{С}$ заливают $m=1\,\text{кг}$ теплой воды. Какова была начальная температура воды, если при установлении теплового равновесия в сосуде расплавилось $560\,\text{г}$ льда? Ответ приведите в градусах цельсия.
}{
\smallskip\noindent\textbf{Решение}\par
При контакте тел с разной температурой, между ними начинается теплообмен. он продолжается до тех пор, пока температуры тел не выровняются. Льда в термосе много, а значит, весь лед не растает, и конечная температура системы будет $0 ^\circ\,\text{С}.$ Поскольку система находится в термосе, теплопотерями можно пренебречь. На плавление льда требуется количество теплоты $Q_1=\Delta m\lambda,$ при охлаждении воды выделяется количество теплоты $Q_2=cm ( t_1 - t_2 ) .$ 
Составим уравнение теплового баланса: 

$Q_1 + Q_2=0.$ 
$\Delta m\lambda + cm ( t_1 - t_2 ) =0.$ 

 В итоге для начальной температуры воды в сосуде с учетом, что $t_1=0 ^\circ C,$ получаем:

\[t_2= \frac{\lambda \Delta m}{cm} = \frac{3{,}3 \cdot 10 ^{ 5 } \cdot 0{,}56}{4200 \cdot 1} =44 ^\circ\,\text{С}.\]

Ответ: $44$ °C.
}

% Source 33074: https://phys-ege.sdamgia.ru/problem?id=33074
\BankTask{33074}{
Лазер испускает световые импульсы с энергией $0{,}1\,\text{Дж}$. Частота повторения импульсов $10\,\text{Гц}$. Через охлаждающую систему лазера прокачивается вода массой $8{,}5\,\text{кг}$, которая в течение часа нагрелась на $10$ \ensuremath{{}^\circ}C. Найдите КПД лазера, если он определяется как отношение излучаемой энергии к потребляемой им.
}{
\smallskip\noindent\textbf{Решение}\par
КПД лазера равен $\eta= \frac{Q_1}{Q_2} ,$ где $Q_{2}$  — количество теплоты, подводимой к лазеру, $Q_{1}$  — количество теплоты, испускаемой лазером за час работы.
Энергия, испускаемая лазером равна суммарной энергии всех фотонов $Q_1=NE$ Число импульсов, испускаемых за указанное время, равно $N=\nu t,$ E  — энергия одного импульса.
$Q_{2}$  — затраченная энергия, равная $Q_2=cm\Delta t + Q_1.$
Таким образом, КПД лазера

\[\eta= \frac{\nu tE}{\nu tE + mc\Delta t} = \frac{10 \cdot 3600 \cdot 0{,}1}{10 \cdot 3600 \cdot 0{,}1 + 4200 \cdot 8{,}5 \cdot 10} \approx 1\%.\]

Ответ: $1$\%.
}

% Source 50827: https://phys-ege.sdamgia.ru/problem?id=50827
\BankTask{50827}{
В калориметре находится $1\,\text{кг}$ льда при температуре −$5$ °C. Какую массу воды, имеющей температуру $20$ °C, нужно добавить в калориметр, чтобы температура его содержимого после установления теплового равновесия оказалась −$2$ °C? Теплообменом с окружающей средой и теплоёмкостью калориметра пренебречь.
}{
\smallskip\noindent\textbf{Решение}\par
В результате теплообмена происходят процессы нагревания льда от $t_{1}$ до установившейся температуры t, охлаждение воды от $t_{2}$ до $t_{\,\text{пл}}=0 ^\circ \,\text{С},$ кристаллизация воды и охлаждение образовавшегося льда от $t_{\,\text{пл}}$ до установившейся температуры t. Для каждого из этих процессов запишем количество теплоты: 

$Q_1=m_1\,\text{с}_{\,\text{л}} ( t - t_1 ) ;$
$Q_2=m_2c_{\,\text{в}} ( t_{\,\text{пл}} - t_2 ) = - m_2c_{\,\text{в}}t_2;$
$Q_3= - \lambda m_2;$
$Q_4=m_2c_{\,\text{л}} ( t - t_{\,\text{пл}} ) =m_2c_{\,\text{л}}t.$

Так как при теплообмене отсутствует теплообмен с окружающей средой, то уравнение теплового баланса имеет вид:

$Q_1 + Q_2 + Q_3 + Q_4=0,$
$m_1\,\text{с}_{\,\text{л}} ( t - t_1 ) - m_2c_{\,\text{в}}t_2 - \lambda m_2 + m_2c_{\,\text{л}}t=0,$

откуда находим массу воды:

\[m_2= \frac{m_1\,\text{с}_{\,\text{л}} ( t - t_1 )}{c_{\,\text{в}}t_2 + \lambda - c_{\,\text{л}}t} = \frac{1 \cdot 2100 \cdot ( - 2 + 5 )}{4200 \cdot 20 + 330000 + 2100 \cdot 2} \approx 0{,}015\,\text{кг}=15\,\text{г}.\]

Ответ: $15\,\text{г}$.
}

% Source 2988: https://phys-ege.sdamgia.ru/problem?id=2988
\BankTask{2988}{
В калориметре находился лед при температуре $t_1= - 5 ^\circ\,\text{С}.$ Какой была масса $m_1$ льда, если после добавления в калориметр $m_2=4\,\text{кг}$ воды, имеющей температуру $t_2=20 ^\circ\,\text{С},$ и установления теплового равновесия температура содержимого калориметра оказалась равной $t=0 ^\circ\,\text{С},$ причем в калориметре была только вода?
}{
\smallskip\noindent\textbf{Решение}\par
Количество теплоты, полученное при нагреве льда, находящегося в калориметре, до температуры $t=0 ^\circ C$: $Q_1=c_1m_1 ( 0 - t_1 )$ ($1$).
Количество теплоты, полученное льдом при его таянии при $0 ^\circ\,\text{С}:$ $Q_2=\lambda m_1$ ($2$).
Количество теплоты, отданное водой при охлаждении ее до $0 ^\circ\,\text{С}:$ $Q=c_2m_2 ( t_2 - 0 )$ ($3$). 
Уравнение теплового баланса: $Q=Q_1 + Q_2$ ($4$). 
Объединяя ($1$)–($4$), получаем: 

\[m_1= \frac{m_2}{c} _2 ( t_2 - 0 ) c_1 ( 0 - t_1 ) + \lambda = \frac{4 \cdot 4200 \cdot 20}{2100 \cdot 5 + 3{,}3 \cdot 10 ^{5}} \approx 1\,\text{кг}.\]

Ответ: $m_1\approx 1\,\text{кг}.$
}

\end{document}
